\chapter{Kết Luận Và Hướng Phát Triển}

\section{Kết luận}
Đề tài bài tập lớn môn Kỹ thuật Vi xử lý về việc thiết kế và chế tạo nguyên mẫu \textbf{"Trạm quan trắc và cảnh báo chất lượng không khí cục bộ sử dụng vi điều khiển ESP32"} đã được nhóm thực hiện nghiên cứu, chế tạo và thử nghiệm thực tế thành công. Hệ thống đã đáp ứng đầy đủ các yêu cầu chức năng, phi chức năng và chỉ tiêu kỹ thuật đề ra ban đầu.

Các kết quả chính đạt được của đề tài bao gồm:
\begin{enumerate}
    \item \textbf{Về phần cứng}: Thiết kế và gia công thành công bo mạch PCB tích hợp đầy đủ mạch bảo vệ, các bộ đệm transistor điều khiển còi báo, LED và quạt hút. Chế tạo buồng đo lấy mẫu khí chủ động bằng nhựa acrylic hoạt động ổn định giúp cải thiện đáng kể độ tin cậy của phép đo cảm biến.
    \item \textbf{Về firmware}: Xây dựng thành công kiến trúc đa tác vụ trên hệ điều hành FreeRTOS chạy song song trên lõi kép của ESP32. Sử dụng hiệu quả các khóa Mutex bảo vệ an toàn đa luồng cho tài nguyên bộ nhớ và bus vật lý I2C. Thiết kế cơ chế offline cache trên SPIFFS hoạt động tin cậy, chống mất mát dữ liệu khi mất kết nối mạng.
    \item \textbf{Về dịch vụ đám mây}: Tích hợp cơ sở dữ liệu Supabase, xây dựng luồng Webhook và Edge Functions TypeScript chạy trên Deno độc lập để gửi thông báo cảnh báo về Telegram Bot tự động mà không cần lưu token bảo mật trên thiết bị nhúng.
    \item \textbf{Về phân tích số liệu thực nghiệm}: Thực hiện thành công đợt thực nghiệm đo đạc kéo dài 55 ngày (1.320 giờ đo liên tục) thu về cơ sở dữ liệu lớn. Huấn luyện các mô hình học máy (Hồi quy tuyến tính và Random Forest) để hiệu chuẩn thành công sai số do độ ẩm của cảm biến bụi Sharp thô bám sát trạm tham chiếu quốc gia.
\end{enumerate}

Nguyên mẫu trạm đo IoT chế tạo có chi phí rất thấp (dưới 2 triệu VNĐ), dễ dàng nhân rộng lắp đặt và đóng vai trò là một mô hình thực hành nhúng chuẩn mực cho sinh viên.

\section{Hạn chế của hệ thống}
Mặc dù đạt được nhiều kết quả tích cực, nguyên mẫu hiện tại của hệ thống vẫn tồn tại một số điểm hạn chế cần khắc phục:
\begin{itemize}
    \item \textbf{Độ bền và trôi cảm biến (Sensor Drift)}: Các cảm biến MQ135 và GP2Y1010AU0F là dòng cảm biến giá rẻ, lớp nhạy cảm bề mặt dễ bị lão hóa sau 6 - 12 tháng sử dụng liên tục, dẫn đến hiện tượng trôi điểm không (offset drift). GP2Y cũng dễ bị bụi bám bẩn vĩnh viễn lên thấu kính nếu buồng đo không được vệ sinh định kỳ.
    \item \textbf{Tính tuyến tính của cảm biến bụi}: Cảm biến bụi Sharp sử dụng quang học LED hồng ngoại, độ nhạy đối với các hạt bụi mịn có đường kính nhỏ hơn $2.5\mu m$ (PM2.5) kém hơn nhiều so với các cảm biến sử dụng chùm tia laser hiện đại, dẫn đến độ phân giải phép đo còn hạn chế.
    \item \textbf{Độ thẩm mỹ vỏ hộp}: Vỏ hộp gia công thủ công bằng nhựa acrylic ghép nối tuy đáp ứng tốt chức năng thông gió buồng đo nhưng kích thước còn khá cồng kềnh, độ thẩm mỹ chưa cao để thương mại hóa thương mại.
\end{itemize}

\section{Hướng phát triển tương lai}
Để nâng cao độ chính xác, tính ứng dụng thực tiễn của hệ thống, nhóm đề xuất một số hướng phát triển và nâng cấp trong tương lai như sau:
\begin{enumerate}
    \item \textbf{Nâng cấp cảm biến bụi laser}: Thay thế cảm biến Sharp GP2Y bằng các cảm biến đo bụi mịn sử dụng công nghệ tán xạ laser hiện đại như PMS5003 hoặc PMS7003. Các cảm biến này tích hợp sẵn bộ vi xử lý lọc nhiễu bên trong, cho phép phân tách và đọc chính xác nồng độ bụi mịn PM1.0, PM2.5 và PM10 một cách độc lập.
    \item \textbf{Nhúng thuật toán hiệu chuẩn trực tiếp vào MCU (TinyML)}: Thay vì chạy mô hình Random Forest hoặc Hồi quy tuyến tính trên máy tính ngoại vi bằng Python, nhóm sẽ nghiên cứu tối ưu hóa và nạp trực tiếp mô hình hiệu chuẩn đã huấn luyện xuống ESP32 sử dụng thư viện TensorFlow Lite cho Vi điều khiển (TinyML). ESP32 sẽ tự động hiệu chuẩn thời gian thực giá trị đo bụi/khí dựa trên nhiệt độ và độ ẩm đọc được trước khi hiển thị và gửi lên cloud.
    \item \textbf{Thiết kế Web Dashboard hoàn chỉnh}: Phát triển ứng dụng Web Front-end (sử dụng Next.js hoặc React) kết nối thời gian thực với cơ sở dữ liệu Supabase qua Supabase Realtime JS để vẽ các biểu đồ phân tích trực quan sinh động cho người dùng.
    \item \textbf{Tối ưu hóa nguồn điện dự phòng}: Thiết kế lại mạch nguồn sạc pin Li-ion tích hợp chip quản lý sạc chuyên dụng (như TP4056 hoặc MP2636) hỗ trợ cắm thêm tấm pin năng lượng mặt trời nhỏ để thiết bị có thể tự cấp nguồn hoạt động độc lập hoàn toàn ngoài trời mà không cần cắm điện AC lưới.
    \item \textbf{Thiết kế vỏ hộp in 3D công nghiệp}: Sử dụng phần mềm CAD (như SolidWorks) để thiết kế vỏ hộp bảo vệ đạt tiêu chuẩn kháng nước bụi IP54, chế tạo bằng công nghệ in 3D nhựa chất lượng cao giúp tăng tính thẩm mỹ và độ bền cơ học cho thiết bị khi đặt ngoài trời.
\end{enumerate}
